% !TEX root = main.tex
\section*{Abstract — 국문 요약}
원료를 사용한 생산과 기업 간 이전이 진행된 뒤에도, 해당 원료에 관한 증빙에서 문제가 발견될 수 있다.
이때 과거에 어떤 원료를 사용했는지 확인하고, 문제가 있는 원료에서 파생된 원장 객체의 추가 소비를 제한하는 대응이 필요할 수 있다.
그러나 입력과 출력의 연결을 숨기는 원장에서는 공개된 검증 결과만으로 후속 대상과 그 소비를 차단할 식별자를 알아내기 어렵다.

본 논문은 비공개 공급망 회계, 승인된 감사자의 이력 추적, 원장 객체의 추가 소비 차단을 연결하는 zkDPP를 제안한다.
거래 증명은 암호화된 부모 참조와 출력의 미래 소비 식별자가 실제 입력·출력에 대응하는지 검증한다.
신뢰된 감사자는 이 정보를 복호화해 기준 시점에 아직 소비되지 않은 후속 대상 전체를 찾고, 해당 대상의 상태를 변경해 추가 소비를 차단한다.

명시한 구성 조건 아래 정확성, 건전성, 프라이버시를 정의하고 증명한다.
정확성은 정상 조건에서 거래·추적·동결이 기대한 결과를 내는 성질이다.
건전성은 잘못된 거래의 수락, 부정확한 추적 결과의 성공 판정, 불완전한 차단의 완료 보고를 배제하고, 차단 상태가 유지되는 동안 소비가 거부되는 성질을 다룬다.
프라이버시는 기록 암호와 감사키 보호 암호의 기밀성을 바탕으로 무작위 오라클 모형에서 분석하며, 참여자의 비밀과 객체의 개방값을 얻지 못한 수동 외부 관찰자에 한정한다.

동결 전에 대상이 소비되거나 대상별 상태를 서로 다른 시점에 확인하는 사례를 통해, 개별 동결 성공과 전체 차단 완료가 다름을 보인다.
차단은 원장 객체의 추가 소비에만 적용되며, 이미 발급된 DPP Claim의 취소나 실물 제품의 회수는 포함하지 않는다.
기호 모델의 사례 검사는 이러한 모델상의 동작을 확인하는 별도 근거이며, 구체적인 암호 방식이 가정을 충족하는지, 실제 구현이 명세에 부합하는지, 성능이 어떠한지는 검증하지 않았다.
